% SGA 3，报告 VII_A，中文重构——第 1 轮。
% 范围：权威底本本地第 29--31 页，4.3--4.3.3，止于本地第 31 页第 5 节之前。
% 合并阅读本页码： 482--484. 再版本印刷页码： 471--473.
% 本段原版页边定位： 442--444.
% 权威底本：Polo--Gille Exp7A-13oct24.pdf,
% SHA-256 D46DB8C79D0385643B6EED4B96900613248375F9B73B4202478A4F57FD3E8522.
% 参照了 Jacob Reinhold 的 CC BY 4.0 Markdown 作为比对谱系。

\medskip
\noindent\textbf{4.3.}
\phantomsection
\label{sga3:VIIA:4.3}

% 权威底本： Exp7A 本地页 29; 合并阅读本页 482; 印刷页 471.
最后，设 \(G\) 是一个取值于阿贝尔群的 \(S\)-群概形。由乘法定义的复合态射
\[
  \mu(G):
  U^p(G)\lhook\joinrel\longrightarrow G_S^p
  \longrightarrow G,
\]
经由 \(V^p(G)\) 分解。因此有态射
\[
  \nu(G):V^p(G)\longrightarrow G
\]
满足
\[
  \nu(G)\circ q'(G)=\mu(G).
\]
所得交换图为
\begin{center}
  \phantomsection
  \label{sga3:VIIA:4.3-verschiebung-master-diagram}
  % SGA3-VII-DIAG-046
  \begin{tikzpicture}[
  baseline=(current bounding box.center),
  >=Stealth,
  every node/.style={font=\normalsize}
]
  \node (gleft) at (-3.3,2.1) {$G$};
  \node (upg)   at (0,2.1) {$\mathrm U^p(G)$};
  \node (gright) at (3.3,2.1) {$G$};
  \node (vpg)   at (-1.65,0.35) {$\mathrm V^p(G)$};
  \node (gbr)   at (1.65,0.35) {$G^{[p/S]}$};
  \node (gps)   at (0,-1.65) {$G^{(p/S)}$};

  \draw[->] (upg) -- node[above] {$\mu(G)$} (gleft);
  \draw[->] (gright) -- node[above] {$\delta(G)$} (upg);
  \draw[->] (vpg) -- node[pos=.52,above,fill=white,inner sep=1pt]
    {$\nu(G)$} (gleft);
  \draw[->] (upg) -- node[pos=.53,right,fill=white,inner sep=1pt]
    {$q'(G)$} (vpg);
  \draw[->] (gbr) -- node[below] {$i(G)$} (vpg);
  \draw[->] (gright) -- node[pos=.5,left,fill=white,inner sep=1pt]
    {$F^{[p]}(G/S)$} (gbr);
  \draw[->] (gbr) -- node[pos=.55,anchor=west,xshift=2pt,fill=white,inner sep=1pt]
    {$\varphi(G)$} (gps);
  \draw[->] (gright) to[out=-92,in=18]
    node[right] {$\operatorname{Fr}(G/S)$} (gps);
  \draw[->,dotted] (gps) to[out=172,in=-87]
    node[left] {$\operatorname{Ver}(G/S)$} (gleft);
  \node at (3.55,-1.65) {$.$};
\end{tikzpicture}

\end{center}

当 \(G\) 在 \(S\) 上平坦时，\(\varphi(G)\) 是同构，于是可定义称为 \emph{Verschiebung} 的态射
\[
  \operatorname{Ver}(G/S):
  G^{(p/S)}
  \longrightarrow
  G
\]
其定义为
\[
  \operatorname{Ver}(G/S)
  =
  \nu(G)\circ i(G)\circ\varphi(G)^{-1}.
\]
当 \(G\) 遍历平坦交换 \(S\)-群概形时，对应
\[
  \operatorname{Ver}:G\longmapsto\operatorname{Ver}(G/S)
\]
显然是函子性同态。因此，\(\operatorname{Ver}(G/S)\) 是群同态。

对每个 \(S\)-概形 \(T\)，复合
\[
  G(T)
  \xrightarrow{\delta(G)(T)}
  U^p(G)(T)
  \xrightarrow{\mu(G)(T)}
  G(T)
\]
将 \(x\in G(T)\) 映到 \(p\cdot x\)。因此可用 \(p\cdot\operatorname{id}_G\) 表示 \(\mu(G)\circ\delta(G)\)，得到经典公式
\[
  \tag{*}\label{sga3:vii:paragraph:4-3:equation:star}
  \operatorname{Ver}(G/S)\circ\operatorname{Fr}(G/S)
  =
  p\cdot\operatorname{id}_G.
\]

\medskip
\noindent\textbf{例 4.3.1.}
\phantomsection
\label{sga3:VIIA:4.3.1}

\smallskip
\noindent\textit{a)}
若 \(G\) 是取值于阿贝尔群的常值 \(S\)-概形，则 \(\operatorname{Fr}(G/S)\) 可等同于 \(G\) 的恒等态射（参见 \SGARefLink{sga3:VIIA:4.1.1.1}{4.1.1.1}）。故
\[
  \operatorname{Ver}(G/S)
  =
  p\cdot\operatorname{id}_G.
\]

\smallskip
\noindent\textit{b)}
若 \(G\) 是类型为 \(M\) 的可对角化 \(S\)-群，则由 \SGARefLink{sga3:VIIA:4.1.4}{4.1.4}~\textit{a)} 有 \(\operatorname{Fr}(G/S)=p\cdot\operatorname{id}_G\)。立即得到 \(\operatorname{Ver}(G/S)\) 是 \(G\) 的恒等态射。

\smallskip
\noindent\textit{c)}
若 \(\mathcal E\) 是平坦 \(\mathcal O_S\)-模，且 \(G=\mathcal V(\mathcal E)\)，则 \(\operatorname{Ver}(G/S)\) 为零，\(p\cdot\operatorname{id}_G\) 也为零。在报告 VII\textsubscript{B} 中，我们将看到域 \(k\) 上的交换代数群 \(G\) 是“幺幂的”，当且仅当存在某个 \(n\)，使复合同态
\begin{center}
  \phantomsection
  \label{sga3:VIIA:4.3.1-iterated-verschiebung-chain}
  % SGA3-VII-DIAG-047
  \(\displaystyle
G^{(p^n)}
  \xrightarrow{\operatorname{Ver}(G^{(p^{n-1})}/S)}
G^{(p^{n-1})}
  \longrightarrow \cdots \longrightarrow
G^{(p)}
  \xrightarrow{\operatorname{Ver}(G/S)}
G
\)

\end{center}
为零。这里
\[
  G^{(p^n)}
  =
  \bigl(G^{(p^{n-1})}\bigr)^{(p)}
\]
如 \SGARefLink{sga3:VIIA:4.1.3}{4.1.3} 中一样。\footnote{此判别准则在报告 VII\textsubscript{B} 中没有明确陈述；见 \textup{[DG70]}，IV，\S3，命题~4.11。}

\medskip
\noindent\textbf{4.3.2.}
\phantomsection
\label{sga3:VIIA:4.3.2}

% 权威底本： Exp7A 本地页 30; 合并阅读本页 483; 印刷页 472.
由于 \(\operatorname{Ver}:G\mapsto\operatorname{Ver}(G/S)\) 是平坦交换 \(S\)-群概形范畴上的函子性同态，方块图
\begin{center}
  \phantomsection
  \label{sga3:VIIA:4.3.2-frobenius-verschiebung-square}
  % SGA3-VII-DIAG-048
  \begin{tikzcd}[row sep=3.3em,column sep=5.2em]
G^{(p)} \arrow[r,"\operatorname{Ver}(G/S)"]
  \arrow[d,swap,"\operatorname{Fr}(G/S)^{(p)}"]
  & G \arrow[d,"\operatorname{Fr}(G/S)"] \\
G^{(p^2)} \arrow[r,"\operatorname{Ver}(G^{(p)}/S)"]
  & G^{(p)}
\end{tikzcd}

\end{center}
交换。这里 \(\operatorname{Fr}(G/S)^{(p)}\) 表示 \(\operatorname{Fr}(G/S)\) 沿基变换 \(\operatorname{fr}(S)\) 的拉回。由 \SGARefLink{sga3:VIIA:4.1.1}{4.1.1}，
\[
  \operatorname{Fr}(G/S)^{(p)}
  =
  \operatorname{Fr}(G^{(p)}/S).
\]
将 \SGARefLink{sga3:VIIA:4.3}{4.3} 的公式 \((*)\) 应用于 \(G^{(p)}\)，得到
\[
  \tag{**}\label{sga3:vii:paragraph:4-3-2:equation:double-star}
  \operatorname{Fr}(G/S)\circ\operatorname{Ver}(G/S)
  =
  \operatorname{Ver}(G^{(p)}/S)
    \circ\operatorname{Fr}(G^{(p)}/S)
  =
  p\cdot\operatorname{id}_{G^{(p)}}.
\]

\medskip
\noindent\textbf{4.3.3.}
\phantomsection
\label{sga3:VIIA:4.3.3}

最后假设 \(G\) 是一个有限局部自由的交换 \(S\)-群。设 \(\mathcal A\) 为其仿射 \(\mathcal O_S\)-代数，\(\pi\) 为将截面 \(x\) 映到 \(x^p\) 的环层 \(\mathcal O_S\) 的自同态。\footnote{编者改变了叙述顺序，先引入下图中出现的对象。}
用 \(\Sigma^p\mathcal A\) 表示 \(\bigotimes_{\mathcal O_S}^p\mathcal A\) 中由对称群作用下不变的截面组成的子代数，并用
\[
  i(\mathcal A):
  \Sigma^p\mathcal A
  \longrightarrow
  \bigotimes_{\mathcal O_S}^p\mathcal A
\]
表示其包含态射。

设
\[
  \Delta^p(\mathcal A):
  \mathcal A
  \longrightarrow
  \bigotimes_{\mathcal O_S}^p\mathcal A
\]
为迭代余代数 \(\mathcal A\) 的对角态射所得的态射。它对应于乘法态射 \(U^p(G)=G_S^p\to G\)。由 \SGARefLink{sga3:VIIA:4.3}{4.3} 开头的讨论，\(\Delta^p(\mathcal A)\) 经由 \(\Sigma^p\mathcal A\) 分解；即它诱导态射
\[
  a(\mathcal A):
  \mathcal A\longrightarrow\Sigma^p\mathcal A
\]
满足
\[
  i(\mathcal A)\circ a(\mathcal A)
  =
  \Delta^p(\mathcal A).
\]

另一方面，设 \(S^p(\mathcal A)\) 是 \(\mathcal A\) 的对称代数中的 \(p\) 次齐次分量，设
\[
  q(\mathcal A):
  \bigotimes_{\mathcal O_S}^p\mathcal A
  \longrightarrow
  S^p(\mathcal A)
\]
为典范投影。乘法
\[
  m^p(\mathcal A):
  \bigotimes_{\mathcal O_S}^p\mathcal A
  \longrightarrow
  \mathcal A
\]
经由 \(S^p(\mathcal A)\) 分解；等价地，它诱导映射
\[
  b(\mathcal A):
  S^p(\mathcal A)\longrightarrow\mathcal A
\]
满足
\[
  b(\mathcal A)\circ q(\mathcal A)
  =
  m^p(\mathcal A).
\]

由于 \(\Sigma^p\mathcal A\) 是 \(V^p(G)\) 的仿射代数，再由 \SGARefLink{sga3:VIIA:4.3}{4.3} 开头的讨论，复合 \(i(G)\circ\varphi(G)^{-1}\) 诱导代数同态
\[
  r(\mathcal A):
  \Sigma^p\mathcal A
  \longrightarrow
  \mathcal A\otimes_\pi\mathcal O_S.
\]
它在形如下式的截面上为零
\[
  \sum_{\sigma\in\mathfrak S_p}
    a_{\sigma(1)}\otimes\cdots\otimes a_{\sigma(p)}
\]
并将 \(a\otimes\cdots\otimes a\) 映到 \(a\otimes_\pi1\)。类似地，\(j(\mathcal A)\) 是 \(\mathcal O_S\)-模态射
\[
  a\otimes_\pi1
  \longmapsto
  q(a\otimes\cdots\otimes a).
\]

% 权威底本： Exp7A 本地页 31; 合并阅读本页 484; 印刷页 473.
这些映射组成交换图
\begin{center}
  \phantomsection
  \label{sga3:VIIA:4.3.3-cartier-duality-algebra-diagram}
  % SGA3-VII-DIAG-049
  \begin{tikzpicture}[
  baseline=(current bounding box.center),
  >=Stealth,
  every node/.style={font=\normalsize}
]
  \node (atl) at (-3.2,2.0) {$\mathcal A$};
  \node (atm) at (0,2.0) {$\displaystyle\bigotimes_{\mathcal O_S}^{p}\mathcal A$};
  \node (atr) at (3.2,2.0) {$\mathcal A$};
  \node (aml) at (-1.65,0.15) {$\Sigma^p\mathcal A$};
  \node (amr) at (1.65,0.15) {$\mathrm S^p(\mathcal A)$};
  \node (ab)  at (0,-1.75) {$\mathcal A\otimes_{\pi}\mathcal O_S$};
  \node at (-4.85,0.15) {$({\mathcal A})$};

  \draw[->] (atl) -- node[above] {$\Delta^p(\mathcal A)$} (atm);
  \draw[->] (atm) -- node[above] {$m^p(\mathcal A)$} (atr);
  \draw[->] (atl) -- node[midway,left,fill=white,inner sep=1pt]
    {$a(\mathcal A)$} (aml);
  \draw[->] (aml) -- node[midway,left,fill=white,inner sep=1pt]
    {$i(\mathcal A)$} (atm);
  \draw[->] (atm) -- node[midway,right,fill=white,inner sep=1pt]
    {$q(\mathcal A)$} (amr);
  \draw[->] (amr) -- node[midway,right,fill=white,inner sep=1pt]
    {$b(\mathcal A)$} (atr);
  \draw[->] (aml) -- node[midway,left,fill=white,inner sep=1pt]
    {$r(\mathcal A)$} (ab);
  \draw[->] (ab) -- node[midway,right,fill=white,inner sep=1pt]
    {$j(\mathcal A)$} (amr);
  \node at (3.55,-1.75) {$.$};
\end{tikzpicture}

\end{center}
复合 \(r(\mathcal A)\circ a(\mathcal A)\) 对应于 Verschiebung 态射 \(\operatorname{Ver}(G/S)\)，而 \(b(\mathcal A)\circ j(\mathcal A)\) 对应于 Frobenius 态射 \(\operatorname{Fr}(G/S)\)。

前述交换图 \((\mathcal A)\) 是自对偶的。事实上，设 \(D\) 为将每个 \(\mathcal O_S\)-模 \(\mathcal M\) 映到其对偶的函子
\[
  D\mathcal M
  =
  \operatorname{Hom}_{\mathcal O_S}
    (\mathcal M,\mathcal O_S).
\]
图 \((\mathcal A)\) 在 \(D\) 下的像恰好是图 \((D\mathcal A)\)。态射
\[
  Dr(\mathcal A),\quad
  Da(\mathcal A),\quad
  Dj(\mathcal A),\quad
  Db(\mathcal A)
\]
分别可等同于
\[
  j(D\mathcal A),\quad
  b(D\mathcal A),\quad
  r(D\mathcal A),\quad
  a(D\mathcal A).
\]
由 \SGARefLink{sga3:VIIA:3.3.1}{3.3.1} 得到：

\begin{quote}
\textit{在有限局部自由的交换 \(S\)-群范畴中，Cartier 对偶互换 Frobenius 态射与 Verschiebung。}
\end{quote}
\footnote{另见 \textup{[DG70]}, IV, \S3, 4.9.}

% 下一定位： Exp7A 本地页 31, 合并阅读本页 484,
% 再版本印刷页 473, 在第 5 节。
